\begin{thebibliography}{99}

\bibitem[Pinotsis \& Miller(2022)]{pinotsis2022}
 Pinotsis, D. A., \& Miller, E. K. (2022). Beyond dimension reduction: Stable electric fields emerge from and allow representational drift. \textit{NeuroImage}, 253, 119058.

\bibitem[Tononi et al.(2016)]{tononi2016}
 Tononi, G., Boly, M., Massimini, M., \& Koch, C. (2016). Integrated information theory: from consciousness to its physical substrate. \textit{Nature Reviews Neuroscience}, 17(7), 450-461.

\bibitem[Singer \& Gray(1995)]{singer1995}
 Singer, W., \& Gray, C. M. (1995). Visual feature integration and the temporal correlation hypothesis. \textit{Annual Review of Neuroscience}, 18, 555-586.

\bibitem[Fries(2015)]{fries2015}
 Fries, P. (2015). Rhythms for cognition: communication through coherence. \textit{Neuron}, 88(1), 220-235.

\bibitem[Shadlen \& Movshon(1999)]{shadlen1999}
 Shadlen, M. N., \& Movshon, J. A. (1999). Synchrony unbound: a critical evaluation of the temporal binding hypothesis. \textit{Neuron}, 24(1), 67-77.

\bibitem[Tsuchiya et al.(2015)]{tsuchiya2015}
 Tsuchiya, N., Wilke, M., Fr\"{a}ssle, S., \& Lamme, V. A. F. (2015). No-report paradigms: extracting the true neural correlates of consciousness. \textit{Trends in Cognitive Sciences}, 19(12), 757-770.

\bibitem[Metzinger(2003)]{metzinger2003}
 Metzinger, T. (2003). \textit{Being No One: The Self-Model Theory of Subjectivity}. Cambridge, MA: MIT Press.

\bibitem[Seth(2013)]{seth2013}
 Seth, A. K. (2013). Interoceptive inference, emotion, and the embodied self. \textit{Trends in Cognitive Sciences}, 17(11), 565-573.

\bibitem[Butlin et al.(2023)]{butlin2023}
 Butlin, P., Long, R., Elmoznino, E., Bengio, Y., Birch, J., et al. (2023). Consciousness in artificial intelligence: insights from the science of consciousness. \textit{arXiv preprint} arXiv:2308.08708.

\bibitem[Bobadilla-Suarez et al.(2020)]{bobadillasuarez2020}
 Bobadilla-Suarez, S., Ahlheim, C., Mehrotra, A., Panos, A., \& Love, B. C. (2020). Measures of neural similarity. \textit{Computational Brain \& Behavior}, 3, 369-383.

\bibitem[DiCarlo \& Cox(2007)]{dicarlo2007}
 DiCarlo, J. J., \& Cox, D. D. (2007). Untangling invariant object recognition. \textit{Trends in Cognitive Sciences}, 11(8), 333-341.

\bibitem[Kriegeskorte \& Diedrichsen(2019)]{kriegeskorte2019}
 Kriegeskorte, N., \& Diedrichsen, J. (2019). Peeling the onion of brain representations. \textit{Annual Review of Neuroscience}, 42, 407-432.

\bibitem[Gallagher(2000)]{gallagher2000}
 Gallagher, S. (2000). Philosophical conceptions of the self: implications for cognitive science. \textit{Trends in Cognitive Sciences}, 4(1), 14-21.

\bibitem[Baars(1997)]{baars1997}
 Baars, B. J. (1997). In the Theater of Consciousness: The Workspace of the Mind. \textit{Oxford University Press}.

\bibitem[Dehaene et al.(2014)]{dehaene2014}
 Dehaene, S., Charles, L., King, J. R., \& Marti, S. (2014). Toward a computational theory of conscious processing. \textit{Current Opinion in Neurobiology}, 25, 76-84.

\bibitem[Friston(2010)]{friston2010}
 Friston, K. (2010). The free-energy principle: a unified brain theory?. \textit{Nature Reviews Neuroscience}, 11(2), 127-138.

\bibitem[Milinkovic \& Aru(2026)]{milinkovic2026}
 Milinkovic, B., \& Aru, J. (2026). On biological and artificial consciousness: A case for biological computationalism. \textit{Neuroscience \& Biobehavioral Reviews}, 181, 106524. \href{https://doi.org/10.1016/j.neubiorev.2025.106524}{doi:10.1016/j.neubiorev.2025.106524}.

\bibitem[Aru et al.(2023)]{aru2023}
 Aru, J., Larkum, M. E., \& Shine, J. M. (2023). The feasibility of artificial consciousness through the lens of neuroscience. \textit{Trends in Neurosciences}, 46(12), 1008--1017. \href{https://doi.org/10.1016/j.tins.2023.09.009}{doi:10.1016/j.tins.2023.09.009}.

\bibitem[Seth(2026a)]{seth2026}
 Seth, A. K. (2026). Conscious artificial intelligence and biological naturalism. \textit{Behavioral and Brain Sciences}, 49, e315. \href{https://doi.org/10.1017/S0140525X25000032}{doi:10.1017/S0140525X25000032}.

\bibitem[Seth(2026b)]{seth2026response}
 Seth, A. K. (2026). The stuff matters: Consciousness, computation, and biology. \textit{Behavioral and Brain Sciences}, 49, e366. \href{https://doi.org/10.1017/S0140525X26106906}{doi:10.1017/S0140525X26106906}.

\bibitem[Gurnee et al.(2026)]{gurnee2026}
 Gurnee, W., et al. (2026). Verbalizable Representations Form a Global Workspace in Language Models. \textit{Transformer Circuits Thread}. \href{https://arxiv.org/abs/2607.15495}{arXiv:2607.15495}.

\bibitem[Aguilera et al.(2022)]{aguilera2022}
 Aguilera, M., Millidge, B., Tschantz, A., \& Buckley, C. L. (2022). How particular is the physics of the free energy principle?. \textit{Physics of Life Reviews}, 40, 24-50.

\bibitem[Biehl et al.(2021)]{biehl2021}
 Biehl, M., Pollock, F. A., \& Kanai, R. (2021). A technical critique of some parts of the free energy principle. \textit{Entropy}, 23(3), 293.

\bibitem[Bruineberg et al.(2022)]{bruineberg2022}
 Bruineberg, J., Dolega, K., Dewhurst, J., \& Baltieri, M. (2022). The Emperor's New Markov Blankets. \textit{Behavioral and Brain Sciences}, 45, e183.

\bibitem[Pinotsis \& Miller(2023a)]{pinotsis2023ephaptic}
 Pinotsis, D. A., \& Miller, E. K. (2023). In vivo ephaptic coupling allows memory network formation. \textit{Cerebral Cortex}, 33(17), 9877-9895.

\bibitem[Pinotsis \& Miller(2023b)]{pinotsis2023cytoelectric}
 Pinotsis, D. A., Fridman, G., \& Miller, E. K. (2023). Cytoelectric coupling: Electric fields sculpt neural activity and "tune" the brain's infrastructure. \textit{Progress in Neurobiology}, 226, 102465.

\bibitem[McFadden(2002)]{mcfadden2002}
 McFadden, J. (2002). The conscious electromagnetic information (cemi) field theory: the hard problem made easy?. \textit{Journal of Consciousness Studies}, 9(8), 45-60.

\bibitem[Pockett(2012)]{pockett2012}
 Pockett, S. (2012). The electromagnetic field theory of consciousness: a testable hypothesis about the characteristics of conscious as opposed to non-conscious fields. \textit{Journal of Consciousness Studies}, 19(11-12), 191-223.

\bibitem[McFadden(2020)]{mcfadden2020}
 McFadden, J. (2020). Integrating information in the brain's EM field: the cemi field theory of consciousness. \textit{Neuroscience of Consciousness}, 2020(1), niaa016.

\bibitem[Fr\"{o}hlich \& McCormick(2010)]{frohlich2010}
 Fr\"{o}hlich, F., \& McCormick, D. A. (2010). Endogenous electric fields may guide neocortical network activity. \textit{Neuron}, 67(1), 129-143.

\bibitem[Anastassiou \& Koch(2015)]{anastassiou2015}
 Anastassiou, C. A., \& Koch, C. (2015). Ephaptic coupling to endogenous electric field activity: why bother?. \textit{Current Opinion in Neurobiology}, 31, 95-103.

\bibitem[Pockett(2002)]{pockett2002}
 Pockett, S. (2002). Difficulties with the electromagnetic field theory of consciousness. \textit{Journal of Consciousness Studies}, 9(4), 51-56.

\bibitem[V\"or\"oslakos et al.(2018)]{voroslakos2018}
 V\"or\"oslakos, M., et al. (2018). Direct effects of transcranial electric stimulation on brain circuits in rats and humans. \textit{Nature Communications}, 9, 483.

\bibitem[Verkhratsky \& Nedergaard(2018)]{verkhratsky2018}
 Verkhratsky, A., \& Nedergaard, M. (2018). Physiology of Astroglia. \textit{Physiological Reviews}, 98(1), 239-389.

\bibitem[Mountcastle(1997)]{mountcastle1997}
 Mountcastle, V. B. (1997). The columnar organization of the neocortex. \textit{Brain}, 120(4), 701-722. \href{https://doi.org/10.1093/brain/120.4.701}{doi:10.1093/brain/120.4.701}.

\bibitem[Horton \& Adams(2005)]{horton2005}
 Horton, J. C., \& Adams, D. L. (2005). The cortical column: a structure without a function. \textit{Philosophical Transactions of the Royal Society B}, 360(1456), 837-862. \href{https://doi.org/10.1098/rstb.2005.1623}{doi:10.1098/rstb.2005.1623}.

\bibitem[Rakic(2008)]{rakic2008}
 Rakic, P. (2008). Confusing cortical columns. \textit{Proceedings of the National Academy of Sciences}, 105(34), 12099-12100. \href{https://doi.org/10.1073/pnas.0807271105}{doi:10.1073/pnas.0807271105}.

\bibitem[Bastos et al.(2021)]{bastos2021}
 Bastos, A. M., Donoghue, J. A., Brincat, S. L., Mahnke, M., Yanar, J., Correa, J., ... \& Miller, E. K. (2021). Neural effects of propofol-induced unconsciousness and its reversal using thalamic stimulation. \textit{eLife}, 10, e60824.

\bibitem[Bajwa et al.(2025)]{bajwa2025}
   Bajwa, I. J., Nilsen, A. S., Skukies, R., Aamodt, A., Ernst, G., Storm, J. F., \& Juel, B. E. (2025). A repeated awakening study exploring the capacity of complexity measures to capture dreaming during propofol sedation. \textit{Scientific Reports}, 15, 32746. \href{https://doi.org/10.1038/s41598-025-12695-z}{doi:10.1038/s41598-025-12695-z}.

\bibitem[Xie et al.(2013)]{xie2013}
 Xie, L., Kang, H., Xu, Q., Chen, M. J., Liao, Y., Thiyagarajan, M., ... \& Nedergaard, M. (2013). Sleep drives metabolite clearance from the adult brain. \textit{Science}, 342(6156), 373-377.

\bibitem[Sykov\'{a} \& Nicholson(2008)]{sykova2008}
 Sykov\'{a}, E., \& Nicholson, C. (2008). Diffusion in brain extracellular space. \textit{Physiological Reviews}, 88(4), 1277-1340.

\bibitem[Barbour(2017)]{barbour2017}
 Barbour, B. (2017). Analysis of claims that the brain extracellular impedance is high and non-resistive. \textit{Biophysical Journal}, 113(7), 1636-1638.

\bibitem[Logothetis et al.(2007)]{logothetis2007}
 Logothetis, N. K., Kayser, C., \& Oeltermann, A. (2007). In vivo measurement of cortical impedance spectrum in monkeys: implications for signal propagation. \textit{Neuron}, 55(5), 809--823. \href{https://doi.org/10.1016/j.neuron.2007.07.027}{doi:10.1016/j.neuron.2007.07.027}.

\bibitem[Gabriel et al.(1996)]{gabriel1996}
 Gabriel, S., Lau, R. W., \& Gabriel, C. (1996). The dielectric properties of biological tissues: II. Measurements in the frequency range 10 Hz to 20 GHz. \textit{Physics in Medicine and Biology}, 41(11), 2251--2269. \href{https://doi.org/10.1088/0031-9155/41/11/002}{doi:10.1088/0031-9155/41/11/002}.

\bibitem[Norton(2005)]{norton2005}
 Norton, J. D. (2005). Eaters of the lotus: Landauer's principle and the return of Maxwell's demon. \textit{Studies in History and Philosophy of Science Part B: Studies in History and Philosophy of Modern Physics}, 36(2), 375-411.

\bibitem[Landauer(1961)]{landauer1961}
 Landauer, R. (1961). Irreversibility and heat generation in the computing process. \textit{IBM Journal of Research and Development}, 5(3), 183-191.

\bibitem[Kibble(1976)]{kibble1976}
 Kibble, T. W. B. (1976). Topology of cosmic domains and strings. \textit{Journal of Physics A: Mathematical and General}, 9(8), 1387-1398.

\bibitem[Sakaguchi(1988)]{sakaguchi1988}
 Sakaguchi, H. (1988). Cooperative phenomena in coupled oscillator systems under external fields. \textit{Progress of Theoretical Physics}, 79(1), 39-46.

\bibitem[Strogatz \& Mirollo(1991)]{strogatz1991}
 Strogatz, S. H., \& Mirollo, R. E. (1991). Stability of incoherence in a population of coupled oscillators. \textit{Journal of Statistical Physics}, 63, 613--635.

\bibitem[Acebr\'{o}n et al.(2005)]{acebron2005}
 Acebr\'{o}n, J. A., Bonilla, L. L., P\'{e}rez Vicente, C. J., Ritort, F., \& Spigler, R. (2005). The Kuramoto model: A simple paradigm for synchronization phenomena. \textit{Reviews of Modern Physics}, 77, 137--185. \href{https://doi.org/10.1103/RevModPhys.77.137}{doi:10.1103/RevModPhys.77.137}.

\bibitem[Ott \& Antonsen(2008)]{ott2008}
 Ott, E., \& Antonsen, T. M. (2008). Low dimensional behavior of large systems of globally coupled oscillators. \textit{Chaos}, 18, 037113. \href{https://doi.org/10.1063/1.2930766}{doi:10.1063/1.2930766}.

\bibitem[D\"orfler \& Bullo(2014)]{dorfler2014}
 D\"orfler, F., \& Bullo, F. (2014). Synchronization in complex networks of phase oscillators: A survey. \textit{Automatica}, 50, 1539--1564. \href{https://doi.org/10.1016/j.automatica.2014.04.012}{doi:10.1016/j.automatica.2014.04.012}.

\bibitem[Breakspear(2017)]{breakspear2017}
 Breakspear, M. (2017). Dynamic models of large-scale brain activity. \textit{Nature Neuroscience}, 20, 340--352. \href{https://doi.org/10.1038/nn.4497}{doi:10.1038/nn.4497}.

\bibitem[Still et al.(2012)]{still2012}
 Still, S., Sivak, D. A., Bell, A. J., \& Crooks, G. E. (2012). Thermodynamics of prediction. \textit{Physical Review Letters}, 109(12), 120604.

\bibitem[Kawai et al.(2007)]{kawai2007}
 Kawai, R., Parrondo, J. M. R., \& Van den Broeck, C. (2007). Dissipation: The phase-space perspective. \textit{Physical Review Letters}, 98(8), 080602.

\bibitem[Horowitz and Esposito(2014)]{horowitz2014}
 Horowitz, J. M., \& Esposito, M. (2014). Thermodynamics with continuous information flow. \textit{Physical Review X}, 4, 031015. \href{https://doi.org/10.1103/PhysRevX.4.031015}{doi:\allowbreak10.1103/\allowbreak PhysRevX.4.\allowbreak031015}.

\bibitem[Ito and Sagawa(2013)]{ito2013}
 Ito, S., \& Sagawa, T. (2013). Information thermodynamics on causal networks. \textit{Physical Review Letters}, 111, 180603. \href{https://doi.org/10.1103/PhysRevLett.111.180603}{doi:\allowbreak10.1103/\allowbreak PhysRevLett.111.\allowbreak180603}.

\bibitem[Corberi et al.(2026)]{corberi2026}
 Corberi, F., dello Russo, S., Messuti, G., Scarpetta, S., \& Smaldone, L. (2026). Thermodynamic learning. \textit{arXiv preprint}. \href{https://arxiv.org/abs/2609.04732}{arXiv:2609.04732}.

\bibitem[Banach(1922)]{banach1922}
 Banach, S. (1922). Sur les op\'{e}rations dans les ensembles abstraits et leur application aux \'{e}quations int\'{e}grales. \textit{Fundamenta Mathematicae}, 3(1), 133-181.

\bibitem[Shenker(2000)]{shenker2000}
  Shenker, O. R. (2000). Logic and entropy. \textit{PhilSci-Archive preprint}. \href{https://philsci-archive.pitt.edu/115/}{philsci-archive.pitt.edu/115}.

\bibitem[Rosas et al.(2026)]{rosas2026}
  Rosas, F. E., Mediano, P. A. M., Luppi, A. I., Boyd, A. B., Berta, M., Jensen, H. J., Gastpar, M., \& Carroll, S. M. (2026). Four faces of information in natural systems. \textit{Nature Reviews Physics}. \href{https://doi.org/10.1038/s42254-026-00971-4}{doi:10.1038/s42254-026-00971-4}.

\bibitem[Sznitman(1991)]{sznitman1991}
  Sznitman, A.-S. (1991). Topics in propagation of chaos. In \textit{\'Ecole d'\'Et\'e de Probabilit\'es de Saint-Flour XIX---1989}, Lecture Notes in Mathematics, vol.~1464, pp.~165--251. Springer. \href{https://doi.org/10.1007/BFb0085169}{doi:10.1007/BFb0085169}.

\bibitem[Dai Pra \& den Hollander(1996)]{daipra1996}
 Dai Pra, P., \& den Hollander, F. (1996). McKean--Vlasov limit for interacting random processes in random media. \textit{Journal of Statistical Physics}, 84, 735--772. \href{https://doi.org/10.1007/BF02179656}{doi:10.1007/BF02179656}.

\bibitem[Bertini et al.(2010)]{bertini2010}
 Bertini, L., Giacomin, G., \& Pakdaman, K. (2010). Dynamical aspects of mean field plane rotators and the Kuramoto model. \textit{Journal of Statistical Physics}, 138, 270--290. \href{https://doi.org/10.1007/s10955-009-9908-9}{doi:10.1007/s10955-009-9908-9}.

\bibitem[Berglund \& Gentz(2002)]{berglund2002}
  Berglund, N., \& Gentz, B. (2002). Pathwise description of dynamic pitchfork bifurcations with additive noise. \textit{Probability Theory and Related Fields}, 122, 341--388. \href{https://doi.org/10.1007/s004400100174}{doi:10.1007/s004400100174}.

\bibitem[Abramsky \& Brandenburger(2011)]{abramsky2011}
  Abramsky, S., \& Brandenburger, A. (2011). The sheaf-theoretic structure of non-locality and contextuality. \textit{New Journal of Physics}, 13, 113036. \href{https://doi.org/10.1088/1367-2630/13/11/113036}{doi:10.1088/1367-2630/13/11/113036}.

\bibitem[Townsend et al.(2015)]{townsend2015}
 Townsend, R. G., Solomon, S. S., Chen, S. C., Pietersen, A. N. J., Martin, P. R., Solomon, S. G., \& Gong, P. (2015). Emergence of complex wave patterns in primate cerebral cortex. \textit{The Journal of Neuroscience}, 35(11), 4657-4662. \href{https://doi.org/10.1523/JNEUROSCI.4509-14.2015}{doi:10.1523/JNEUROSCI.4509-14.2015}.

\bibitem[Xu et al.(2023)]{xu2023}
 Xu, Y., Long, X., Feng, J., \& Gong, P. (2023). Interacting spiral wave patterns underlie complex brain dynamics and are related to cognitive processing. \textit{Nature Human Behaviour}, 7(7), 1196-1215. \href{https://doi.org/10.1038/s41562-023-01626-5}{doi:10.1038/s41562-023-01626-5}.

\bibitem[Nicoletti et al.(2026)]{nicoletti2026}
 Nicoletti, G., Di Terlizzi, I., \& Busiello, D. M. (2026). Balancing information and dissipation with partially observed fluctuating signals. \textit{Physical Review Letters}, 137(12), 127101. \href{https://doi.org/10.1103/8gx2-rbns}{doi:10.1103/8gx2-rbns}.

\bibitem[Kiyooka et al.(2026)]{kiyooka2026}
 Kiyooka, D., Oomoto, I., Kitazono, J., Saito, Y., Kobayashi, M., Matsubara, C., Kobayashi, K., Murayama, M., \& Oizumi, M. (2026). Single-cell resolution functional networks during unconsciousness are segregated into spatially intermixed modules. \textit{Cell Reports}, 45(2), 116902. \href{https://doi.org/10.1016/j.celrep.2025.116902}{doi:10.1016/j.celrep.2025.116902}.

\bibitem[Oomoto et al.(2026)]{oomoto2026}
 Oomoto, I., Kiyooka, D., Oizumi, M., \& Murayama, M. (2026). Multi-area single-cell calcium imaging dataset of the mouse cortex across wakefulness, sleep, and anesthesia. \textit{Scientific Data}. \href{https://doi.org/10.1038/s41597-026-08202-2}{doi:10.1038/s41597-026-08202-2}.

\bibitem[Attwell \& Laughlin(2001)]{attwell2001}
 Attwell, D., \& Laughlin, S. B. (2001). An energy budget for signaling in the grey matter of the brain. \textit{Journal of Cerebral Blood Flow \& Metabolism}, 21(10), 1133-1145. \href{https://doi.org/10.1097/00004647-200110000-00001}{doi:10.1097/00004647-200110000-00001}.

\bibitem[Harris et al.(2012)]{harris2012}
 Harris, J. J., Jolivet, R., \& Attwell, D. (2012). Synaptic energy use and supply. \textit{Neuron}, 75(5), 762-777. \href{https://doi.org/10.1016/j.neuron.2012.08.019}{doi:10.1016/j.neuron.2012.08.019}.

\bibitem[Nicholls \& Ferguson(2013)]{nicholls2013}
 Nicholls, D. G., \& Ferguson, S. J. (2013). \textit{Bioenergetics 4}. Academic Press.


\bibitem[Baer et al.(1989)]{baer1989}
 Baer, S. M., Erneux, T., \& Rinzel, J. (1989). The slow passage through a Hopf bifurcation: delay, memory effects, and resonance. \textit{SIAM Journal on Applied Mathematics}, 49(1), 55--71. \href{https://doi.org/10.1137/0149003}{doi:10.1137/0149003}.

\bibitem[Friedman et al.(2010)]{friedman2010}
 Friedman, E. B., Sun, Y., Moore, J. T., Hung, H.-T., Meng, Q. C., Perera, P., Joiner, W. J., Thomas, S. A., Eckenhoff, R. G., Sehgal, A., \& Kelz, M. B. (2010). A conserved behavioral state barrier impedes transitions between anesthetic-induced unconsciousness and wakefulness: evidence for neural inertia. \textit{PLoS ONE}, 5(7), e11903. \href{https://doi.org/10.1371/journal.pone.0011903}{doi:10.1371/journal.pone.0011903}.

\bibitem[Hudson et al.(2014)]{hudson2014}
 Hudson, A. E., Calderon, D. P., Pfaff, D. W., \& Proekt, A. (2014). Recovery of consciousness is mediated by a network of discrete metastable activity states. \textit{Proceedings of the National Academy of Sciences}, 111(25), 9283--9288. \href{https://doi.org/10.1073/pnas.1408296111}{doi:10.1073/pnas.1408296111}.

\bibitem[Proekt \& Hudson(2018)]{proekt2018}
 Proekt, A., \& Hudson, A. E. (2018). A stochastic basis for neural inertia in emergence from general anaesthesia. \textit{British Journal of Anaesthesia}, 121(1), 86--94. \href{https://doi.org/10.1016/j.bja.2018.02.035}{doi:10.1016/j.bja.2018.02.035}.
\end{thebibliography}
